\section*{开始之前：把空间、极限、集合分开理解}
本讲的核心问题是：一列点越来越接近某处时，那个地方是否仍在原集合里？这句话会依次变成极限、闭集、紧集，并为下一讲的连续性服务。全文在有限维欧氏空间讨论；维数记为 $N$，数列编号记为 $k$，不要把二者混淆。$x^k$ 表示第 $k$ 个向量，$x^k_j$ 表示它的第 $j$ 个坐标，通常不是乘方。

约定 $B_r(x)=\{y:\norm{y-x}<r\}$ 为\emph{开球}，$\overline B_r(x)=\{y:\norm{y-x}\le r\}$ 为闭球。原课件对 $N,B$ 的使用前后不完全一致，讲解统一采用这个约定。$E'$ 是聚点集，$\overline E$ 是闭包，$E^\circ$ 是内部，$\partial E$ 是边界，$E^c$ 是相对于当前空间的补集。撇号在此不是求导。

\paragraph{证明的三种基本动作}证明 $A\subseteq B$：任取 $x\in A$，推出 $x\in B$。证明等号：做两次包含。证明``任意 $\eps>0$''的结论：先让别人任意给出一个正数 $\eps$，然后依题意选择半径或编号；选择不能依赖尚未给出的变量。

\Slide{1}{本讲主题：有限维空间中的极限}
本页是课件封面。本讲围绕 $\R^N$ 中的点和点集展开，与 Debreu《Theory of Value》第一章第6节对应。学习目标不是记住所有图形，而是能用``邻域、数列、补集''三种语言表达同一件事。证明考试中，必须写清楚所取的点、半径、数列分别属于哪里。

\Slide{2}{本讲在整门课中的位置}
第2节集合与第3节函数提供语言，第5节实数提供上确界和完备性，本节把实数结论推广到有限个坐标。后面连续函数需要闭集和紧集，连续对应需要序列与闭图，不动点定理需要非空、紧、凸。因此本节优先掌握：极限的量词、闭集的序列判别、紧性的有限覆盖性质，以及连通与区间的关系。

\Slide{3}{一个点就是一串有顺序的实数}
$\R^N=\underbrace{\R\times\cdots\times\R}_{N\text{ 个}}$。例如 $(2,-1,4)\in\R^3$，其第2坐标是 $-1$。坐标有顺序，$(1,2)$ 与 $(2,1)$ 不同。经济学里一个商品组合可以记成 $x=(x_1,\ldots,x_N)$，但本讲暂不要求坐标非负。证明中写 $x=y$ 等价于逐一证明 $x_j=y_j$ 对所有 $j$ 成立。

\Slide{4}{二维坐标的几何读法}
$\R^2$ 中第一坐标控制左右，第二坐标控制上下。点 $(a,b)$ 是直线 $x_1=a$ 与 $x_2=b$ 的交点。区分点 $(a,b)$ 与区间 $(a,b)\subset\R$：同一括号的意思由上下文决定。以后讨论一个平面集合是否开，必须允许向\emph{所有平面方向}作小移动，不能只沿一条线移动。

\Slide{5}{三维坐标与二维切面}
$\R^3$ 中一个等式 $x_1=a$ 通常确定一个平面，三个独立坐标等式才确定一个点。这样理解比想象高维图形更可靠：高维证明基本上都是逐坐标计算。例如 $x=(1,0,-2)$，其第三坐标是 $-2$；保持其他坐标不动、只改变第三坐标，仍在 $\R^3$ 中。

\Slide{6}{坐标子空间与维数}
给定指标集合 $I\subseteq\{1,\ldots,N\}$，只允许 $j\in I$ 的坐标自由变化，其余坐标固定为0，得到坐标子空间，维数为 $|I|$。例如在 $\R^3$ 中令 $I=\{1,3\}$，得到 $\{(a,0,c):a,c\in\R\}$，是二维平面。原点 $0=(0,\ldots,0)$ 不是空集；一个是点，一个是没有元素的集合。

\Slide{7}{三种距离与一般的 $p$ 距离}
令 $a_j=|x_j-y_j|$，则 $\rho_1(x,y)=\sum a_j$，$\rho_2(x,y)=\sqrt{\sum a_j^2}$，$\rho_\infty(x,y)=\max_j a_j$。从 $(0,0)$ 到 $(3,4)$，三者分别为7、5、4。$p$ 距离是 $(\sum a_j^p)^{1/p}$；$p\ge1$ 才有课件要求的三角不等式。所谓 $\ell^p$ 是距离的名称，通常只有 $p=2$ 被称为欧氏距离。后面的开闭性在这些距离下相同，但距离数值并不相同。

\Slide{8}{为什么 $p\to\infty$ 会得到最大坐标}
课件用对数和洛必达法则计算；含有 $a_j=0$ 时 $\log a_j$ 需要额外处理。以下夹逼证明更短且覆盖所有情形。令 $M=\max_j a_j$。因为某个 $a_j=M$，而每个 $a_j\le M$，有
\[
 M^p\le\sum_{j=1}^N a_j^p\le NM^p,
 \qquad M\le\rho_p(x,y)\le N^{1/p}M.
\]
$N^{1/p}\to1$，故夹逼得 $\rho_p\to M=\rho_\infty$。若 $M=0$，左右本来全为0。证明模板是``最大项给下界，项数乘最大项给上界''。

\Slide{9}{距离公理与距离函数的连续性}
非负性与对称性直接来自绝对值；三角不等式表示绕经第三点不会更短。由两次三角不等式得
\[
 |\rho(x,y)-\rho(u,v)|\le\rho(x,u)+\rho(y,v).
\]
具体地，$\rho(x,y)\le\rho(x,u)+\rho(u,v)+\rho(v,y)$ 给一个方向，再交换两对点给反方向。如果 $x_k\to x,y_k\to y$，右边趋0，于是 $\rho(x_k,y_k)\to\rho(x,y)$。这证明不依赖尚未学过的连续函数运算法则，是后面使用距离连续性的依据。一般 $p$ 距离的三角不等式推导见第12页讲解。

\Slide{10}{开球、闭球、去心邻域}
开球要求距离严格小于 $r$；闭球允许等于 $r$；去心球 $B_r(x)\setminus\{x\}$ 额外排除中心。半径始终要求 $r>0$。例如在实数线上，开球是 $(x-r,x+r)$，闭球是 $[x-r,x+r]$。在最大范数下球是方块，在欧氏范数下是圆或球。去心条件用于聚点：集合只有中心一个点时，不能因此称中心为聚点。

\Slide{11}{范数是到原点的距离}
$\norm{x}_p=\rho_p(x,0)$ 是向量的大小，$\rho_p(x,y)=\norm{x-y}_p$ 是两个点的距离。$\norm{x}_\infty=\max_j|x_j|$，不能漏掉绝对值。例如 $(-5,2)$ 的最大范数是5，不能取2。范数的输入是一个向量，距离的输入是两个点；这一区分能避免写错极限条件。

\Slide{12}{范数比较与三角不等式的依据}
\paragraph{比较不等式的证明}设 $1\le p\le q<\infty$。若 $x\ne0$，令 $b_j=|x_j|/\norm{x}_p$，则 $0\le b_j\le1$ 且 $\sum b_j^p=1$。于是 $\sum b_j^q\le1$，得到 $\norm{x}_q\le\norm{x}_p$。另一方面，对凹函数 $t^{p/q}$ 的平均值不等式给出
\[
 \frac1N\sum |x_j|^p\le\left(\frac1N\sum |x_j|^q\right)^{p/q},
 \quad\norm{x}_p\le N^{1/p-1/q}\norm{x}_q.
\]
$q=\infty$ 时直接用第8页的最大项估计。$x=0$ 时显然成立。
\paragraph{三条范数公理}非负且仅零向量范数为0；$\norm{\alpha x}=|\alpha|\norm{x}$；$\norm{x+y}\le\norm{x}+\norm{y}$。对 $p=1,\infty$，逐坐标使用 $|x_j+y_j|\le|x_j|+|y_j|$ 即可。对 $p=2$，先由 $\sum_j(x_j-ty_j)^2\ge0$ 的判别式推出柯西不等式 $|\sum x_jy_j|\le\norm{x}_2\norm{y}_2$，再展开 $\norm{x+y}_2^2$ 即得结论。
\paragraph{一般 $p>1$ 的补充证明}令 $q=p/(p-1)$。Young 不等式 $ab\le a^p/p+b^q/q$ 可由函数 $a^p/p-ab+b^q/q$ 对 $a\ge0$ 的最小值为0证明。将两个向量分别除以其 $p$、$q$ 范数，逐项应用并求和，就得 Hölder 不等式 $\sum|u_jv_j|\le\norm{u}_p\norm{v}_q$。记 $S=\sum|x_j+y_j|^p$，则
\[
 S\le\sum(|x_j|+|y_j|)|x_j+y_j|^{p-1}
 \le(\norm{x}_p+\norm{y}_p)S^{1/q}.
\]
$S>0$ 时除以 $S^{1/q}$ 得三角不等式；$S=0$ 时显然成立。这些补充只用于解释公理的来源，后续可直接使用。

\Slide{13}{有界：一个统一的围栏}
集合 $E$ 有界是 $\exists M<\infty\ \forall x\in E,\norm{x}\le M$。必须先选一个 $M$ 管住整个集合；``对每个点存在它自己的 $M$''对任何集合都成立，不能表达有界。最大范数下 $\norm{x}_\infty\le M$ 等价于每个坐标在 $[-M,M]$ 中，因此得到边长 $2M$ 的闭立方体。\textbf{注意}只有最大范数的球是该立方体；欧氏范数的球不是立方体。有限维范数比较说明选哪一种范数不影响有界性。

\Slide{14}{向量收敛与逐坐标收敛}
$x^k\to x$ 的完整读法是：对任意误差 $\eps>0$，存在整数 $K$，使所有 $k\ge K$ 都满足 $\norm{x^k-x}<\eps$。允许 $K$ 依赖误差，但不能只让某些很晚的项靠近。
\paragraph{等价性证明}向量收敛推出逐坐标收敛，因为 $|x^k_j-x_j|\le\norm{x^k-x}_\infty$。反过来，对每个 $j$ 取 $K_j$ 使误差小于 $\eps$，再令 $K=\max_{1\le j\le N}K_j$，则全部坐标同时小于 $\eps$。一般 $p$ 范数时可先把坐标误差控制为 $\eps/N^{1/p}$。课件中不同范数下同一个 $\eps$ 的不等式并非逐项等价；等价的是收敛性质。

\Slide{15}{极限唯一与收敛必有界}
若 $x^k\to a$ 且 $x^k\to b$，则 $\norm{a-b}\le\norm{a-x^k}+\norm{x^k-b}\to0$，所以 $a=b$。若想完全用量词：假设 $d=\norm{a-b}>0$，取两边误差均小于 $d/3$ 就有 $d<2d/3$，矛盾。
收敛数列一定有界：从某项后 $\norm{x^k-x}<1$，于是尾部范数不超过 $\norm{x}+1$；前面有限项的范数有最大值，把两者取最大即可。反过来有界不一定收敛，如 $(-1)^k$。

\Slide{16}{内点、外点、边界点：量词不能写错}
内点：\emph{存在}一个球全部落在 $E$ 中。外点：\emph{存在}一个球完全不碰 $E$。边界点：\emph{每一个}球既碰到 $E$ 又碰到 $E^c$。\textbf{课件修正：}本页第(3)项写``存在一个混合邻域''不够；内点附近取很大的球也可能混合。例如 $E=(0,1)$，点 $1/2$ 是内点，半径2的球却同时碰到 $E$ 与补集。边界必须经得起任意缩小邻域。

\Slide{17}{聚点、导集、内部、闭包}
聚点 $x\in E'$：每个球中有无穷多个 $E$ 的点。内点一定是聚点，因为非空开球含无穷多个点；聚点不必是内点，例如 $E=\{1/k:k\ge1\}$ 的0。聚点也不必属于集合。内部 $E^\circ$ 收集内点；闭包 $\overline E=E\cup E'$ 将遗漏的极限位置补上。
对上述 $E$，有 $E^\circ=\varnothing$，$E'=\{0\}$，$\overline E=E\cup\{0\}$。不要把闭包与导集混同：前者还保留每一个孤立点。

\Slide{18}{聚点的数列判别：构造互不相同的点}
\paragraph{充分性}若互不相同的 $x_k\in E$ 趋于 $x$，任何球 $B_r(x)$ 都包含全部足够晚的项，因此包含无穷多个不同的点，故 $x\in E'$。
\paragraph{必要性}假设 $x\in E'$。从 $E\cap B_1(x)$ 选 $x_1\ne x$；已选 $x_1,\ldots,x_k$ 后，从 $E\cap B_{1/(k+1)}(x)$ 中排除有限集合 $\{x,x_1,\ldots,x_k\}$，仍有点可选，记为 $x_{k+1}$。于是点互不相同，且 $\norm{x_k-x}<1/k\to0$。
\paragraph{容易漏的条件}若允许数列恒为 $x$，只能说明 $x\in\overline E$，不能说明 $x\in E'$。``互不相同''是本定理区别于闭包判别的关键。

\Slide{19}{导集的单调性与有限并公式}
$A\subseteq B$ 时，任意球中无穷多个 $A$ 点也都是 $B$ 点，所以 $A'\subseteq B'$，进而 $\overline A\subseteq\overline B$。
要证 $(A\cup B)'=A'\cup B'$，容易的一边来自单调性。对另一边用逆否命题更清楚：若 $x\notin A'$ 且 $x\notin B'$，分别有半径 $r,s$ 的球只含有限多个 $A,B$ 点。半径 $\min(r,s)$ 的球只含有限多个 $A\cup B$ 点，故 $x\notin(A\cup B)'$。重复这个论证得到任意\emph{有限}并的公式；无限并不能照搬，例见第43页第3题。

\Slide{20}{内部的单调性与有限交公式}
若 $A\subseteq B$，包含于 $A$ 的小球也包含于 $B$，所以 $A^\circ\subseteq B^\circ$。由此 $(A\cap B)^\circ\subseteq A^\circ\cap B^\circ$。反向取 $x\in A^\circ\cap B^\circ$，存在 $r_A,r_B>0$ 分别使小球落入两集合；令 $r=\min(r_A,r_B)>0$，就有 $B_r(x)\subseteq A\cap B$。因此等号成立。
课件略去的步骤是把两个可能不同的半径合并成一个。有限交可以取有限个正数的最小值；无限多个正数的下确界可能为0，这正是无限交不保开性的原因。

\Slide{21}{Bolzano--Weierstrass：有界无限集必有聚点}
定理同时需要有界与无限。$\Z$ 无限但无界，没有有限位置的聚点；$\{0,1\}$ 有界但有限，也没有聚点。聚点可能落在集合外，例如 $\{1/k\}$ 的聚点0。与数列版本的联系：任意有界数列都有收敛子列；若只出现有限多个不同值，至少一个值出现无限次，取常值子列；若出现无限多个值，使用集合版本和递增下标挑选即可。

\Slide{22}{二分立方体证明：不断保留无限多个点}
把有界集 $E$ 放进 $[-M,M]^N$。沿每个坐标二分，得到 $2^N$ 个闭小方块。若每块只含有限个 $E$ 点，则有限块合起来仍只有有限点，与 $E$ 无限矛盾。因此至少一块仍含无限个点。选这一块再二分，如此构造
\[
 C_0\supseteq C_1\supseteq\cdots,\qquad
 \text{边长}(C_k)=2M/2^k,\quad |C_k\cap E|=\infty.
\]
小方块可以在边界上相交，不影响``有限个有限集之并仍有限''的推理。

\Slide{23}{二分证明的收尾：共同点就是聚点}
每个 $C_k$ 是 $N$ 个闭区间的乘积；逐坐标应用闭区间套定理，得到唯一共同点 $x$。给定 $r>0$，选 $k$ 使 $\sqrt N(2M/2^k)<r$，则 $C_k$ 的欧氏直径小于 $r$。因为 $x\in C_k$，故 $C_k\subset B_r(x)$。而 $C_k\cap E$ 无限，所以每个 $B_r(x)$ 都包含无限个 $E$ 点，$x$ 是聚点。
\textbf{课件修正：}对角线系数应为固定维数的 $\sqrt N$，不能误写成不断变化的子块编号的平方根。最大范数下直径直接等于边长。

\Slide{24}{孤立点与可数性的证明起点}
$x$ 是 $E$ 的孤立点，正确条件是 $x\in E$ 且\emph{存在} $r>0$ 使 $B_r(x)\cap E=\{x\}$，等价于去心球与 $E$ 无交。\textbf{课件修正：}$x\notin D(x,r)$ 总是成立，单独这一句不能定义孤立点。
若所有点都孤立，对每个 $x\in E$ 选 $r_x>0$ 使 $B_{2r_x}(x)\cap E=\{x\}$，再在 $B_{r_x}(x)$ 中选一个有理坐标点 $q_x\in\Q^N$。有理坐标点的存在来自每一坐标上有理数稠密；下一页证明这些 $q_x$ 不会重复。

\Slide{25}{把孤立点注入有理坐标点}
若不同的 $x,y$ 选中了同一个 $q$，则
$\norm{x-y}<r_x+r_y\le2\max(r_x,r_y)$。
若最大者是 $r_x$，则 $y\in B_{2r_x}(x)\cap E$，与该交集只有 $x$ 矛盾；另一情形同理。故 $x\mapsto q_x$ 是单射。$\Q^N$ 是有限个可数集之积，仍可数，所以 $E$ 可数。
本课把 $E\cap E'=\varnothing$ 称为所有点孤立，把 $E'=\varnothing$ 称为 discrete。其他教材常把前者也称为离散子空间，答题应先写条件再写术语。$\{1/k\}$ 的每一点都孤立，但0仍是外部聚点。

\Slide{26}{聚点也可以用去心邻域定义}
等价条件是 $x\in E'\iff\forall r>0,\ (B_r(x)\setminus\{x\})\cap E\ne\varnothing$。正向显然：无限多个点不可能只有中心一个。若课件说一个``含 $x$ 的邻域''却以 $y$ 为中心，只要 $x\in B_r(y)$，仍可找到以 $x$ 为中心的小球 $B_{r-\norm{x-y}}(x)\subset B_r(y)$。因为半径差严格为正，两种邻域表述等价。

\Slide{27}{一个额外点为什么能推出无限多个点}
逆向假设每个去心球都碰到 $E$。若某个 $B_r(x)$ 只含有限个异于 $x$ 的点 $y_1,\ldots,y_m$，取
$0<s<\min(r,\norm{y_1-x},\ldots,\norm{y_m-x})$，则去心球 $B_s(x)\setminus\{x\}$ 不碰 $E$，矛盾。若原球根本没有异点，已经直接矛盾。因此每个球包含无限点。
课件递归选距离严格下降的点也能证明点数无限；但``严格下降的正数列''并不自动趋0，例如 $1+1/k$。若要同时构造收敛列，应额外把第 $k$ 次半径限制到 $1/k$。

\Slide{28}{邻域换中心与内部的局部稳定性}
若 $B_r(x)\subseteq E$，且 $y\in B_r(x)$，记 $d=\norm{x-y}<r$。对 $z\in B_{r-d}(y)$，三角不等式给 $\norm{z-x}<r-d+d=r$，故 $z\in E$。所以原球中的每个点也是内点。这一小计算将反复出现：\emph{可移动的余量等于原半径减去已走过的距离}。若从一个非中心邻域反推中心邻域，也用同样的余量公式。

\Slide{29}{闭集的定义与三个基本例子}
$E$ 闭是指 $E'\subseteq E$，等价于 $\overline E=E$。更适合考试的版本：只要 $x_k\in E$ 且 $x_k\to x$，就有 $x\in E$。证明这个版本时，若极限就是某项值则已经在 $E$；否则可抽出趋向 $x$ 的不同点，故 $x\in E'$。
闭长方体的每个坐标不等式在取极限后保留，所以闭；$\R^N$ 和空集也闭。有理数稠密说明 $\overline\Q=\R$，因此 $\Q$ 不闭。连续函数上下水平集的判别在下一页证明；开球闭包等于闭球要求半径正，见第31页。

\Slide{30}{连续性与上下水平集闭：补上子列步骤}
\paragraph{正向}若 $f$ 连续，且 $x_k\to x$、$f(x_k)\ge t$，则 $f(x)=\lim f(x_k)\ge t$，故 $\{f\ge t\}$ 闭；$\{f\le t\}$ 同理。
\paragraph{反向}若在 $x$ 不连续，则存在 $\eps_0>0$，对每个 $k$ 都能选 $\norm{x_k-x}<1/k$ 而 $|f(x_k)-f(x)|\ge\eps_0$。每项至少属于以下一类：$f(x_k)\le f(x)-\eps_0$ 或 $f(x_k)\ge f(x)+\eps_0$。两类中至少一类有无限多个下标，抽取该类子列。例如第一类得到闭集 $\{f\le f(x)-\eps_0\}$ 中的数列趋于 $x$，但 $x$ 不在该闭集，矛盾。第二类同理。
\textbf{关键补足：}不能假定所有项天然落在同一边；必须先用无限鸽巢原理抽子列。

\Slide{31}{开球的闭包与稠密}
设 $r>0$。若 $x_k\in B_r(a)$ 且 $x_k\to x$，距离连续性给 $\norm{x-a}\le r$，故 $\overline{B_r(a)}\subseteq\overline B_r(a)$。反向对 $x\in\overline B_r(a)$，取 $x_k=a+(1-1/k)(x-a)$，则 $\norm{x_k-a}=(1-1/k)\norm{x-a}<r$，且 $x_k\to x$。因此等号成立。
$A$ 在 $S$ 中稠密，准确表述是 $A\subseteq S$ 且 $S\subseteq\overline A$，或相对闭包 $\overline A^{\,S}=S$。若直接写环境空间中的 $\overline A=S$，还暗中要求 $S$ 闭。后面的第84--86页使用前一种一般定义。

\Slide{32}{抽象闭包运算：四条规则}
课件列出：扩大集合、保持包含、做两次等于一次、空集仍为空集。把运算记为 $c$，就是
$A\subseteq c(A)$，$A\subseteq B\Rightarrow c(A)\subseteq c(B)$，$c(c(A))=c(A)$，$c(\varnothing)=\varnothing$。
这些规则描述抽象闭包算子。\textbf{边界提醒：}仅这四条还不能保证通常拓扑闭包的有限并公式；拓扑闭包还满足 $c(A\cup B)=c(A)\cup c(B)$。本讲的欧氏空间闭包满足更强的这些性质，抽象段落不要无条件套用全部欧氏结论。

\Slide{33}{导集与闭包本身都是闭集}
证明 $E'$ 闭：设 $x\in(E')'$。给定 $r>0$，选 $y\in E'$ 满足 $0<\norm{y-x}<r/2$，再取
$s<\min(r-\norm{y-x},\norm{y-x})$。由于 $y$ 是 $E$ 的聚点，可选 $z\in E\cap B_s(y)$。这个小球不含 $x$，又包含于 $B_r(x)$，因此找到了 $E$ 中异于 $x$ 的近点。由第26--27页，$x\in E'$。
再用有限并导集公式：$(\overline E)'=(E\cup E')'=E'\cup(E')'\subseteq E'\subseteq\overline E$，所以 $\overline E$ 闭。课件略去的小球半径选择正是证明严密性所在。

\Slide{34}{有限个闭集之并仍闭：导集方法}
设 $F=F_1\cup\cdots\cup F_m$，每个 $F_i$ 闭。由第19页的有限并公式，$F'=F_1'\cup\cdots\cup F_m'\subseteq F_1\cup\cdots\cup F_m=F$，于是 $F$ 闭。注意 ``finite closed sets'' 在这里表示\emph{有限个闭集}，不是每个闭集本身只有有限个点。考试写明 $m<\infty$，防止把定理误用到无限并。

\Slide{35}{同一个有限并结论：闭包方法}
本页给另一种代数写法：
\[
 \overline{\bigcup_{i=1}^m F_i}
 =\bigcup_{i=1}^m(F_i\cup F_i')
 =\bigcup_{i=1}^m\overline{F_i}
 =\bigcup_{i=1}^m F_i.
\]
第一个等号用闭包等于本集并导集及有限并导集公式；最后一个等号用每个 $F_i$ 闭。这与上一页是一条证明的两种记号。复习时应能解释每个等号来自哪条定义，不能仅抄下一串等号。

\Slide{36}{可数个闭集之交仍闭：导集方法}
设 $F=\bigcap_{k=1}^\infty F_k$。因为 $F\subseteq F_k$，单调性给 $F'\subseteq F_k'\subseteq F_k$，对\emph{每个} $k$ 都成立，于是 $F'\subseteq\bigcap_k F_k=F$。这里不用交换``求导集''与``无限交''：一般只有包含，不应写成等号。这个证明实际上完全没有用到可数性，下一组页面将推广到任意指标集。

\Slide{37}{可数闭交的闭包写法}
对每个 $k$ 有 $\overline F\subseteq\overline{F_k}=F_k$，所以 $\overline F\subseteq F$。再结合总成立的 $F\subseteq\overline F$，得到 $F=\overline F$。本页比上一页多用一次``双向包含推出相等''的动作。若 $F$ 为空，也一样成立，不需要单独假设交集非空。

\Slide{38}{任意闭集族的交：不需要数得出来}
设 $\Gamma$ 可以不可数，$F=\bigcap_{\gamma\in\Gamma}F_\gamma$。任取 $x\in F'$，因为 $F\subseteq F_\gamma$，可知 $x\in F_\gamma'\subseteq F_\gamma$，对所有 $\gamma$ 成立，所以 $x\in F$。原页公式残留可数指标 $n$ 的位置应统一改为 $\gamma\in\Gamma$。若指标集为空，约定交集为整个环境空间，仍为闭集。

\Slide{39}{任意闭交的闭包证明与笔误}
延续第37页，把 $k$ 换成 $\gamma$：$\overline F\subseteq\overline{F_\gamma}=F_\gamma$，故 $\overline F\subseteq\bigcap_\gamma F_\gamma=F$。\textbf{课件修正：}每个 $F_\gamma$ 闭意味着 $\overline{F_\gamma}=F_\gamma$；不意味着 $\overline{F_\gamma}\subseteq F$，因为交集通常比单个集合小。不要把单调性方向倒过来。

\Slide{40}{为什么无限并会出现新极限点}
$F_k=[1/(k+1),1/k]$ 都闭，但 $\bigcup_kF_k=(0,1]$ 不含聚点0，所以不闭。具体看 $x_k=1/k$：每项落在某个闭区间，但没有一个固定的区间包含无限多个趋0的项。
一般总有
\[
 \bigcup_\gamma\overline{E_\gamma}\subseteq\overline{\bigcup_\gamma E_\gamma},
 \qquad \overline{\bigcap_\gamma E_\gamma}\subseteq\bigcap_\gamma\overline{E_\gamma}.
\]
第一式对每个 $E_\gamma\subseteq\bigcup E_\gamma$ 用单调性再取并；第二式对 $\bigcap E_\gamma\subseteq E_\gamma$ 用单调性再取交。等号一般不成立。

\Slide{41}{闭包是包含原集合的最小闭集}
令 $\mathcal F=\{F:F\text{ 闭且 }A\subseteq F\}$，$C=\bigcap_{F\in\mathcal F}F$。由于 $\overline A$ 本身闭且包含 $A$，它是被求交的成员之一，故 $C\subseteq\overline A$。另一方面每个 $F\in\mathcal F$ 满足 $\overline A\subseteq\overline F=F$，所以 $\overline A\subseteq C$。两边合并得 $C=\overline A$。
``最小''按集合包含比较，不是按元素数量比较；$\mathcal F$ 至少含有整个空间，因此并非无从取交。

\Slide{42}{递减非空紧集的交非空：简洁完整证明}
设 $F_1\supseteq F_2\supseteq\cdots$，各集非空、闭且 $F_1$ 有界。从每个 $F_k$ 任取 $x_k\in F_k$。整个数列都在 $F_1$ 中，故有收敛子列 $x_{k_j}\to x$。任意固定 $m$，当 $k_j\ge m$ 时，$x_{k_j}\in F_{k_j}\subseteq F_m$。由 $F_m$ 闭得到 $x\in F_m$。因为 $m$ 任意，$x\in\bigcap_mF_m$。
这个证明自动处理重复集合、重复点，不必按课件分情况。交集未必只有一点；加上直径趋0才唯一。若去掉有界性，$F_k=[k,\infty)$ 给出空交；去掉闭性，$F_k=(0,1/k)$ 也给出空交。

\Slide{43}{闭集习题详解}
1. 不必有限：$\Z$ 的任意子集都无有限聚点，故在 $\R$ 中闭。2. 不成立：$E_1=(0,1),E_2=(1,2)$，交为空而闭包交含1。3. 不成立：$E_k=\{1/k\}$，每个导集为空，但并集的导集含0。4. 闭包是闭的且含 $E$，任何闭集含 $E$ 也含其全部极限点，故恰为交。5. 对每个整数 $m$ 置 $A_m=\{x\in F:f(x)\le m\}$。若 $A_m$ 无限，紧性给一列互异点收敛到 $x_0\in F'$，与沿这些点函数趋于无穷矛盾。故每个 $A_m$ 有限，$F=\bigcup_m A_m$ 可数。6. 若 $(x_n,y_n)\to(x,y)$ 且 $f(x_n)\ge y_n$，连续性和不等式保极限给 $f(x)\ge y$。7. 取 $\Q+n\sqrt2$（$n\in\Z$）：各集合稠密可数，若相交则 $(n-m)\sqrt2\in\Q$，推出 $n=m$。

\Slide{44}{内点与内部}
$x\in E^\circ$ 指存在 $r>0$ 使 $B(x,r)\subseteq E$。内部不是边界。若 $y\in B(x,r)$，取 $s=r-\norm{x-y}>0$，三角不等式给 $B(y,s)\subseteq B(x,r)$。因此球中每点都是内点，$E^\circ$ 开。内部的单调性来自球的包含，幂等性来自内部本身是开集。

\Slide{45}{内点与内部}
本页继续“内点与内部”。完整定义、条件和论证见本讲原页第 44 页下方。

\Slide{46}{内点与内部}
本页继续“内点与内部”。完整定义、条件和论证见本讲原页第 44 页下方。

\Slide{47}{有限开集的交}
任取 $x\in\bigcap_{i=1}^mG_i$，各取 $r_i>0$ 使 $B(x,r_i)\subseteq G_i$。有限多个正数的最小值 $r$ 仍正，故 $B(x,r)\subseteq\bigcap_iG_i$。无限交不能照搬：$\bigcap_{n\ge1}(-1/n,1/n)=\{0\}$ 不开。

\Slide{48}{有限开集的交}
本页继续“有限开集的交”。完整定义、条件和论证见本讲原页第 47 页下方。

\Slide{49}{任意开集的并}
若 $x\in\bigcup_iG_i$，只需选一个含 $x$ 的 $G_j$；由其开性取球 $B(x,r)\subseteq G_j\subseteq\bigcup_iG_i$。这不需要所有 $G_i$ 共同的正半径，所以任意甚至不可数并仍开。与有限交比较，关键是“存在一个”与“对所有”的区别。

\Slide{50}{任意开集的并}
本页继续“任意开集的并”。完整定义、条件和论证见本讲原页第 49 页下方。

\Slide{51}{任意开集的并}
本页继续“任意开集的并”。完整定义、条件和论证见本讲原页第 49 页下方。

\Slide{52}{任意开集的并}
本页继续“任意开集的并”。完整定义、条件和论证见本讲原页第 49 页下方。

\Slide{53}{实开集的区间分解}
对 $a\in G$，把包含 $a$ 且包含于 $G$ 的开区间取并，得极大开区间 $I_a$。两个这样的区间相交，则并仍是区间，因此极大性使其相等。故不同 $I_a$ 不交且覆盖 $G$。每个非空区间含有理数，选其中一个；不同区间对应不同有理数，故只有可数多个区间。端点允许无穷。

\Slide{54}{实开集的区间分解}
本页继续“实开集的区间分解”。完整定义、条件和论证见本讲原页第 53 页下方。

\Slide{55}{欧氏开集的方块分解}
把空间用边长 $2^{-k}$ 的半开方块分割，保留完全包含于 $G$ 且未被更粗已选方块覆盖的方块。每层可数，总层数可数；嵌套网格保证所选方块两两不交。给定 $x\in G$，存在球包含于 $G$，足够细的含 $x$ 方块直径小于该球半径，所以 $x$ 最终被选中方块覆盖。半开约定避免公共边界重复。

\Slide{56}{欧氏开集的方块分解}
本页继续“欧氏开集的方块分解”。完整定义、条件和论证见本讲原页第 55 页下方。

\Slide{57}{可数拓扑基}
所有有理中心、正有理半径的球构成可数族。若 $x\in G$ 开，选 $B(x,r)\subseteq G$，再选有理中心 $q$ 足够近于 $x$，及有理半径 $s$ 满足 $\norm{x-q}<s<r-\norm{x-q}$，便有 $x\in B(q,s)\subseteq G$。因此任何开集都是某些基球的并。

\Slide{58}{开集习题与必要纠错}
1. 原命题假：$F=\{0\}\subset\R^N$ 闭而边界只有一个点。2. $x$ 不在补集闭包等价于存在球避开补集，即 $x\in E^\circ$；由此 $\partial E=\overline E\setminus E^\circ$。3. 当 $x_2\ne0$ 函数连续；$(a,0)$ 且 $a\ne0$ 时沿正弦取 $\pm1$ 的序列得不同极限，在原点则 $|f|\le|x_1|\to0$，故不连续点为 $(\R\setminus\{0\})\times\{0\}$，不闭。4. 直接用内部、边界、闭包的分解。5. 若 $y=x+v$ 且 $\norm v\le r_0$，由 $G$ 开，$x+h\in G$ 对足够小 $h$ 成立，因而 $y+h=(x+h)+v\in A$，故 $A$ 开。6. 对每个与 $F$ 相交的基球选一个 $F$ 中点，所得集可数且在 $F$ 中稠密；$F$ 闭保证闭包等于 $F$。7. 离散集 $E$ 不交其导集 $E'$，且 $E'$ 闭，故 $E=\overline E\cap(\R^N\setminus E')$。8. 先证不减：若存在 $a<b$ 使 $f(a)>f(b)$，对每个 $y\in(f(b),f(a))$，令 $t_y$ 是 $[a,b]$ 中满足 $f(t_y)=y$ 的最后一点。连续性及介值定理保证它存在且在 $(a,b)$，并且在其右侧一直有 $f(t)<y$，否则连续性又给更晚的同水平点。若 $t_y$ 处可导，则右差商非正，故 $f'(t_y)\le0$，所以 $t_y$ 必在可数例外集。不同 $y$ 给不同 $t_y$，这把不可数实区间单射到可数集，矛盾。故 $f$ 不减。若 $f(a)=f(b)$，不减性使整个 $[a,b]$ 上恒定；从区间选一个不在可数例外集中的内点，该点导数为0，与假设正导数矛盾。因此严格递增。本题不必要求例外点也可导，但必须保持全区间连续。

\Slide{59}{紧性与有限子覆盖}
开覆盖是开集族，其并包含目标集合；有限子覆盖必须从原覆盖取有限成员。有限维中闭有界等价于紧。证明闭有界给有限子覆盖：若没有，把含 $E$ 的闭方块不断二分，每次保留不能有限覆盖的一个子方块。得到直径趋零的套方块及 $x_n\in E$，其极限 $x\in E$（用闭性）。某原覆盖开集含 $x$，故含 $B(x,r)$；足够小的整个所选方块都落入该球，可由一个开集覆盖，矛盾。另一证明的 Lebesgue 数步骤：若无统一半径，取 $x_n$ 使 $B(x_n,1/n)$ 不被单个覆盖成员包含；抽子列 $x_n\to x$，取 $B(x,r)\subset U$，当 $\norm{x_n-x}<r/2,1/n<r/2$ 时矛盾。闭性和有界性缺一不可：$(0,1)$ 可被 $(1/n,1)$ 覆盖却不能有限覆盖，$\N$ 的局部小区间覆盖也不能有限覆盖。

\Slide{60}{紧性与有限子覆盖}
本页继续“紧性与有限子覆盖”。完整定义、条件和论证见本讲原页第 59 页下方。

\Slide{61}{紧性与有限子覆盖}
本页继续“紧性与有限子覆盖”。完整定义、条件和论证见本讲原页第 59 页下方。

\Slide{62}{紧性与有限子覆盖}
本页继续“紧性与有限子覆盖”。完整定义、条件和论证见本讲原页第 59 页下方。

\Slide{63}{紧性与有限子覆盖}
本页继续“紧性与有限子覆盖”。完整定义、条件和论证见本讲原页第 59 页下方。

\Slide{64}{紧性与有限子覆盖}
本页继续“紧性与有限子覆盖”。完整定义、条件和论证见本讲原页第 59 页下方。

\Slide{65}{紧集的邻域与逆定理}
若紧 $F\subset G$ 开，给每个 $y\in F$ 选 $B(y,r_y)\subset G$，用半径 $r_y/2$ 的球有限覆盖 $F$，取有限 $r_y/2$ 的最小值 $\delta>0$。任意 $\norm h<\delta$ 都使 $F+h\subset G$。反之紧集必有界，因为球 $B(0,n)$ 覆盖它且有限子覆盖给统一半径；必闭，因为对 $z\notin E$，以 $x\in E$ 为中心取半径 $\norm{x-z}/3$ 的球，有限覆盖后可在 $z$ 周围找到避开 $E$ 的正半径球。

\Slide{66}{紧集的邻域与逆定理}
本页继续“紧集的邻域与逆定理”。完整定义、条件和论证见本讲原页第 65 页下方。

\Slide{67}{紧性习题}
1. 原覆盖的闭包族当然仍覆盖 $E$，因为每个成员都扩大。2. 任意两闭区间相交给出 $a_\alpha\le b_\beta$ 对所有指标成立；故 $s=\sup_\alpha a_\alpha\le b_\beta$，又 $a_\beta\le s$，所以 $s$ 属于全部区间。任取固定区间可保证确界有限。3. 若非空闭集无孤立点，在每个已有小闭区间内可选两个不同的集内点及互不交更小闭区间；递归构造二叉树并令直径趋零。每条无限0--1路径由完备性给一个集内极限，不同路径给不同点，从而不可数，矛盾。4. 对 $\eps>0$，开集 $U_n=\{x:f_n(x)<\eps\}$ 递增并覆盖紧 $F$，有限子覆盖可合为某个 $U_N$。故所有 $n\ge N,x\in F$ 有 $0\le f_n(x)<\eps$，这正是一致收敛（Dini 型证明）。

\Slide{68}{闭包与内部的抽象公理}
闭包算子 $c$ 满足扩张、单调、幂等及 $c(\varnothing)=\varnothing$；定义 $i(A)=S\setminus c(S\setminus A)$。补集反向包含把扩张变为收缩，单调仍为单调，幂等由两次取补消去得到，$i(S)=S$。反向公式 $c(A)=S\setminus i(S\setminus A)$，因此两构造互逆。闭集是闭包不动集：任意闭集交 $F$ 满足 $c(F)\subseteq c(F_j)=F_j$ 对所有 $j$ 成立，加扩张性得相等。若 $A\subset F$ 闭，则 $c(A)\subset F$；且 $c(A)$ 自己闭，所以它是所有含 $A$ 闭集的交。注意这四条是较一般闭包系统公理，未必单独推出有限闭并性质，欧氏拓扑还有相应额外性质。

\Slide{69}{闭包与内部的抽象公理}
本页继续“闭包与内部的抽象公理”。完整定义、条件和论证见本讲原页第 68 页下方。

\Slide{70}{闭包与内部的抽象公理}
本页继续“闭包与内部的抽象公理”。完整定义、条件和论证见本讲原页第 68 页下方。

\Slide{71}{闭包与内部的抽象公理}
本页继续“闭包与内部的抽象公理”。完整定义、条件和论证见本讲原页第 68 页下方。

\Slide{72}{闭包与内部的抽象公理}
本页继续“闭包与内部的抽象公理”。完整定义、条件和论证见本讲原页第 68 页下方。

\Slide{73}{闭包与内部的抽象公理}
本页继续“闭包与内部的抽象公理”。完整定义、条件和论证见本讲原页第 68 页下方。

\Slide{74}{闭包与内部的抽象公理}
本页继续“闭包与内部的抽象公理”。完整定义、条件和论证见本讲原页第 68 页下方。

\Slide{75}{闭包与内部的抽象公理}
本页继续“闭包与内部的抽象公理”。完整定义、条件和论证见本讲原页第 68 页下方。

\Slide{76}{闭包的序列判据与相对闭性}
$x\in\overline A$ 当且仅对任意 $n$ 可选 $a_n\in A$ 使 $\norm{x-a_n}<1/n$。闭集即任何集内收敛序列的极限仍在集内。闭包幂等可用二段逼近：先在 $\overline A$ 找距 $x$ 小于 $\eps/2$ 的点，再在 $A$ 找距它小于 $\eps/2$ 的点。任意闭交用序列同时属于每个因子；有限闭并则用有限抽屉原理找到无限多项属于某一个闭集的子列。相对闭包为 $S\cap\overline A$，只保留属于背景空间 $S$ 的极限；例如 $(0,1)$ 在自身中闭，却不在 $\R$ 中闭。

\Slide{77}{闭包的序列判据与相对闭性}
本页继续“闭包的序列判据与相对闭性”。完整定义、条件和论证见本讲原页第 76 页下方。

\Slide{78}{闭包的序列判据与相对闭性}
本页继续“闭包的序列判据与相对闭性”。完整定义、条件和论证见本讲原页第 76 页下方。

\Slide{79}{闭包的序列判据与相对闭性}
本页继续“闭包的序列判据与相对闭性”。完整定义、条件和论证见本讲原页第 76 页下方。

\Slide{80}{闭包的序列判据与相对闭性}
本页继续“闭包的序列判据与相对闭性”。完整定义、条件和论证见本讲原页第 76 页下方。

\Slide{81}{闭包的序列判据与相对闭性}
本页继续“闭包的序列判据与相对闭性”。完整定义、条件和论证见本讲原页第 76 页下方。

\Slide{82}{闭包的序列判据与相对闭性}
本页继续“闭包的序列判据与相对闭性”。完整定义、条件和论证见本讲原页第 76 页下方。

\Slide{83}{闭包的序列判据与相对闭性}
本页继续“闭包的序列判据与相对闭性”。完整定义、条件和论证见本讲原页第 76 页下方。

\Slide{84}{任意欧氏子集都可分}
枚举有理中心有理半径的球，对每个与 $S$ 相交的球选择一个点 $s_j\in S$，形成可数集合 $D$。给定 $x\in S,\eps>0$，由可数基选球 $B$ 满足 $x\in B\subset B(x,\eps)$，其选点 $s_j$ 与 $x$ 距离小于 $\eps$。故 $D\subset S\subset\overline D$。不能直接取 $S\cap\Q^N$，它可能为空，例如 $S=\{\sqrt2\}$。

\Slide{85}{任意欧氏子集都可分}
本页继续“任意欧氏子集都可分”。完整定义、条件和论证见本讲原页第 84 页下方。

\Slide{86}{任意欧氏子集都可分}
本页继续“任意欧氏子集都可分”。完整定义、条件和论证见本讲原页第 84 页下方。

\Slide{87}{内部、边界与外部}
内部是 $\R^N\setminus\overline{E^c}$，外部是 $\R^N\setminus\overline E$，边界是 $\overline E\cap\overline{E^c}$。边界点的每个球都同时遇到集合及补集。三者互不交并覆盖空间；证明只需对一点分类：若有球全在 $E$ 则内部，若有球全在补集则外部，否则边界。相对空间中所有补集和闭包都应相对 $S$ 计算。

\Slide{88}{内部、边界与外部}
本页继续“内部、边界与外部”。完整定义、条件和论证见本讲原页第 87 页下方。

\Slide{89}{内部、边界与外部}
本页继续“内部、边界与外部”。完整定义、条件和论证见本讲原页第 87 页下方。

\Slide{90}{内部、边界与外部}
本页继续“内部、边界与外部”。完整定义、条件和论证见本讲原页第 87 页下方。

\Slide{91}{有界、紧与有限积}
有界是存在统一 $M$ 使每个点的范数不大于 $M$。有限积的闭性由坐标收敛和各因子闭性给出；有界性由各因子界的平方和给出。因此有限个紧集的积紧。有限维条件很关键，无限维中闭且有界一般不保证紧。

\Slide{92}{有界、紧与有限积}
本页继续“有界、紧与有限积”。完整定义、条件和论证见本讲原页第 91 页下方。

\Slide{93}{连通的多种判据}
连通指不能分成两个非空不交的相对开集；取补立刻等价于不能分成两个非空不交的相对闭集，也等价于没有非平凡既开又闭子集。$\partial A=\overline A\setminus A^\circ=\varnothing$ 等价于 $A$ 既开又闭，故得边界判据。若不交集合 $A,B$ 覆盖空间且互不碰对方闭包，它们在空间中均闭也均开，给分离判据。最后，非平凡 clopen 集 $A$ 的示性函数 $1_A:X\to\{0,1\}$ 连续，因为每个离散开集的逆像均开；反之非常值连续函数到 $\{0,1\}$ 的两个纤维给分离。实数连通集必是区间：若中间缺点 $c$，两侧相对开集分离它。区间必连通：若分离且 $a$ 在一侧、$b>a$ 在另一侧，取 $s=\sup\{t\in[a,b]:t\text{ 在第一侧}\}$，利用任一侧的相对开性都矛盾。

\Slide{94}{连通的多种判据}
本页继续“连通的多种判据”。完整定义、条件和论证见本讲原页第 93 页下方。

\Slide{95}{连通的多种判据}
本页继续“连通的多种判据”。完整定义、条件和论证见本讲原页第 93 页下方。

\Slide{96}{连通的多种判据}
本页继续“连通的多种判据”。完整定义、条件和论证见本讲原页第 93 页下方。

\Slide{97}{连通的多种判据}
本页继续“连通的多种判据”。完整定义、条件和论证见本讲原页第 93 页下方。

\Slide{98}{连通的多种判据}
本页继续“连通的多种判据”。完整定义、条件和论证见本讲原页第 93 页下方。

\Slide{99}{连通的多种判据}
本页继续“连通的多种判据”。完整定义、条件和论证见本讲原页第 93 页下方。

\Slide{100}{连通的多种判据}
本页继续“连通的多种判据”。完整定义、条件和论证见本讲原页第 93 页下方。
